\section{The line field}\label{sec:linefield}

A \emph{direction field} is a map $\sigma:\Z^4\to\{0,1,2,3\}$; its links
are $F=\{(x,\sigma(x)):x\in\Z^4\}$. At each site put $\pi_x(\mu)=
\sigma(x+e_\mu)$. The plaquette at $x$ in the plane $\{\mu,\nu\}$ contains
\begin{equation}\label{eq:count}
[\sigma(x)\in\{\mu,\nu\}]+[\pi_x(\mu)=\nu]+[\pi_x(\nu)=\mu]
\end{equation}
links of $F$.

\begin{lemma}\label{lem:cycle}
Every plaquette contains exactly one link of $F$ if and only if, at every
site $x$, $\pi_x$ fixes $\sigma(x)$ and permutes the other three directions
cyclically.
\end{lemma}

\begin{proof}
Fix $x$, let $s=\sigma(x)$ and let $R$ be the other three directions. For the
plane $\{s,\nu\}$ the count~\eqref{eq:count} is $1+[\pi_x(s)=\nu]+[\pi_x(\nu)
=s]$, which equals one for all $\nu\in R$ exactly when $\pi_x$ fixes $s$ and
maps $R$ into $R$. For a plane $\{\mu,\nu\}\subset R$ the count is
$[\pi_x(\mu)=\nu]+[\pi_x(\nu)=\mu]$. Draw an arrow $\mu\to\pi_x(\mu)$ whenever
$\pi_x(\mu)\in R\setminus\{\mu\}$; each vertex has at most one outgoing arrow,
and the condition asks for exactly one arrow on each of the three pairs. Three
arrows on three vertices with out-degree at most one give out-degree one
everywhere, and such a tournament is a directed three-cycle.
\end{proof}

From now on $\sigma$ is the field of period two of Table~\ref{tab:sigma} in
Appendix~\ref{app:table}; in each row $\pi_x$ fixes $\sigma(x)$ and is a
three-cycle on the other directions, so every plaquette contains exactly one
link of $F$.\footnote{A field with one link of $F$ at every site and one in
every plaquette can exist only in four dimensions: a site carries
$d(d-1)/2$ plaquettes at which it is the base point and a link lies in
$2(d-1)$ plaquettes, and the two counts agree only for $d=4$.} Call the links
outside $F$ \emph{active} and write $A$ for their set.

\begin{lemma}\label{lem:structure}\leavevmode
\begin{enumerate}[label=\textup{(\roman*)},leftmargin=*]
\item $F$ is a disjoint union of bi-infinite straight lines, one through every
site; hence $F\cap E(\Lambda)$ is a forest for every box $\Lambda$.
\item An active link $e=(x,\mu)$ lies in six plaquettes. In four of them the
link of $F$ shares an end point with $e$; in the other two it is opposite to
$e$.
\item Every active link shares a plaquette with exactly twelve other active
links, two in each of its plaquettes, and two links share at most one
plaquette.
\end{enumerate}
\end{lemma}

\begin{proof}
(i) Since $\pi_x$ fixes $\sigma(x)$, $\sigma(x+e_{\sigma(x)})=\sigma(x)$, and since
$\sigma$ has period two, also $\sigma(x-e_{\sigma(x)})=\sigma(x)$. If $(x-e_\nu,\nu)
\in F$, straightness at $x-e_\nu$ gives $\sigma(x)=\nu$, so the incoming link of
$F$ at $x$ is unique and continues the outgoing one. A line meets a box in a
segment.

(ii) A link of $F$ sharing an end point with $e$ is the outgoing or incoming
link of $F$ at $x$, of direction $s=\sigma(x)\ne\mu$, or at $x+e_\mu$, of
direction $t=\pi_x(\mu)\in R\setminus\{\mu\}$. They lie in the two plaquettes of
$e$ in the plane $\{\mu,s\}$ and the two in the plane $\{\mu,t\}$, four distinct
plaquettes as $s\ne t$. In the two plaquettes of the remaining plane
$\{\mu,\rho\}$ the unique link of $F$ touches neither end of $e$, so it is the
opposite link.

(iii) Each plaquette of $e$ contains $e$, one link of $F$ and two other links.
\end{proof}

